% ============================================================
%  演習用・解答編で共通の紙面設計
%
%  見出しは本文に出さず、柱（左上）と目次にだけ回す。
%  本文に残るのは問題文と、その終わりを示す細い罫線だけにする。
%  原稿 generated/{workbook,answers}/*.tex は
%  tools/split-pdf-sources.mjs が生成する。
%
%  この後で \wbend を定義すると紙面が決まる。
%    演習用 … 罫線から下をそのまま解答欄にする（何も足さない）
%    解答編 … 罫線の下に見通しと解答を続ける
% ============================================================

\usepackage[top=18mm, bottom=14mm, left=20mm, right=20mm, headsep=5mm]{geometry}
\usepackage{fancyhdr}

% 目次は「年度・区分」までにとどめて1ページに収め、
% 問題名は PDF のしおりからだけ辿れるようにする。
\hypersetup{bookmarksdepth=2}

\makeatletter

\newcommand{\wb@year}{}
\newcommand{\wb@sec}{}
\newcommand{\wb@prob}{}
\newcommand{\wb@sep}{\,\textbar\,}

\pagestyle{fancy}
\fancyhf{}
\fancyhead[L]{\normalfont\footnotesize\textcolor{black!60}{%
  \wb@year\wb@sep\wb@sec\wb@sep\textbf{\wb@prob}}}
\fancyhead[R]{\normalfont\footnotesize\textcolor{black!60}{\thepage}}
\renewcommand{\headrulewidth}{0.4pt}
\renewcommand{\footrulewidth}{0pt}
\fancypagestyle{plain}{%
  \fancyhf{}%
  \fancyhead[R]{\normalfont\footnotesize\textcolor{black!60}{\thepage}}%
  \renewcommand{\headrulewidth}{0pt}%
}

% 章・節は柱の材料と目次の見出しにするだけで、本文には何も出さない。
\newcommand{\wbchapter}[1]{%
  \clearpage
  \renewcommand{\wb@year}{#1}%
  \phantomsection\addcontentsline{toc}{chapter}{#1}%
}
\newcommand{\wbsection}[1]{%
  \clearpage
  \renewcommand{\wb@sec}{#1}%
  \phantomsection\addcontentsline{toc}{section}{#1}%
}
% 問題は必ず新しいページから始める。題名は柱と目次へ送る。
\newcommand{\wbproblem}[1]{%
  \clearpage
  \renewcommand{\wb@prob}{#1}%
  \phantomsection\addcontentsline{toc}{subsection}{#1}%
  \par\noindent\ignorespaces
}
% 問題文の終わりを示す細い罫線。
\newcommand{\wbrule}{%
  \par\medskip
  \noindent\textcolor{black!25}{\rule{\linewidth}{0.4pt}}%
  \par
}

\makeatother
